% Section 4.3, from lectures/L04/README.md: the objectives, the questions,
% the further reading, and what the next chapter does with all of it.


\section{Review}
\label{c4:sec:review}

\needspace{6\baselineskip}
\subsection*{What you should be able to do now}
After this chapter, you should be able to:
\begin{itemize}
  \item Compute the loop gain of any of Chapter~\ref{c3:ch:filters}'s configurations, and the gain error that follows.
  \item Say how much open-loop gain a stated accuracy needs, and notice when the answer is unreasonable.
  \item State what feedback does to input impedance, output impedance and distortion, with the direction right in each case.
  \item Convert between gain-bandwidth product, closed-loop gain and closed-loop bandwidth.
  \item Say what a Sallen-Key section buys over two cascaded RC sections with a buffer between them.
  \item Write down the diode equation, and say what a 60 mV change of its voltage does.
  \item Explain why plain Newton-Raphson fails on a diode from a cold start, and what a simulator does about it.
\end{itemize}

\needspace{6\baselineskip}
\subsection*{Questions to test yourself}
You should be able to answer:
\begin{enumerate}
  \item An amplifier has $10^5$ of open-loop gain. You use it at a closed-loop gain of 1000. What is the gain error, and is that amplifier a good choice?
  \item Feedback divides distortion by $1 + T$. What distortion remains, and why can it not be removed by more feedback?
  \item Two amplifiers have the same gain-bandwidth product. One is used at a gain of 10 and one at 100. Which has more loop gain at 10 kHz, and by how much?
  \item Negative feedback becomes positive when the loop phase reaches 180 degrees. How many first-order poles does it take to get there, and what does that say about how many stages a loop can have?
  \item A diode carries 1 mA at 0.65 V. What voltage does it need for 10 mA? For 1 microamp?
  \item Why does a nonlinear solver need a convergence criterion at all, when a linear one just solves?
  \item Your Newton-Raphson loop converges in 6 iterations on one circuit and 400 on another that differs only in a resistor value. What is different about the second one?
\end{enumerate}

\needspace{6\baselineskip}
\subsection*{Further reading}
\begin{itemize}
  \item \emph{The Art of Electronics}, Horowitz and Hill, chapter 4, for feedback, and chapter 1 for the diode.
\end{itemize}

\needspace{6\baselineskip}
\subsection*{Looking ahead}
\begin{itemize}
  \item Chapter~\ref{c5:ch:transistor} begins Part~2. The transistor as a device with an exponential, a current gain and a saturation voltage, and with as little semiconductor physics as can be managed.
\end{itemize}
